\section{Момент силы}

\subsection{Момент силы}

\subsubsection{Определение}
\textbf{Момент силы} (вращательный момент) --- векторная физическая величина, характеризующая вращательное действие силы на тело.

\begin{center}
\begin{tblr}{
  colspec = {X[c]X[l]},
  vlines,
  hlines,
  row{1} = {font=\bfseries,mode=text},
  row{odd} = {bg=gray!10},
  cells = {m},
  rowsep = 6pt
}
Формула & Название \\
$\displaystyle \mathbf{M} = [\mathbf{r}, \mathbf{F}] = \mathbf{r} \times \mathbf{F}$ & \textit{Момент силы относительно точки:} векторное произведение радиус-вектора $\mathbf{r}$ (из точки О в точку приложения силы) на вектор силы $\mathbf{F}$ \\
$\displaystyle M = r \cdot F \cdot \sin \alpha = F \cdot d$ & \textit{Модуль момента силы:} $d$ --- плечо силы (кратчайшее расстояние от точки О до линии действия силы) \\
$\displaystyle F_\perp = F \sin\alpha$ & \textit{Тангенциальная составляющая:} компонента силы, перпендикулярная радиус-вектору \\
$\displaystyle F_r = F \cos\alpha$ & \textit{Радиальная составляющая:} компонента силы, направленная вдоль радиус-вектора \\
$\displaystyle \mathbf{M} = \mathbf{r} \times \mathbf{F}_\perp$ & \textit{Момент силы через тангенциальную составляющую:} радиальная составляющая не создаёт момента \\
\end{tblr}
\end{center}

Направление вектора определяется \textit{по правилу «буравчика» (правой руки)}. Вектор $\mathbf{M}$ перпендикулярен плоскости, в которой лежат $\mathbf{r}$ и $\mathbf{F}$.

\subsubsection{Аддитивность моментов}

Моменты сил, действующих на тело относительно одной и той же точки, складываются векторно. Если на тело действуют несколько сил $\mathbf{F}_1, \mathbf{F}_2, \dots, \mathbf{F}_n$, то полный момент относительно точки $O$ равен векторной сумме их моментов:
\[
\mathbf{M}_{\text{полн}} = \sum_{i=1}^n \mathbf{M}_i = \sum_{i=1}^n [\mathbf{r}_i, \mathbf{F}_i] = \mathbf{r}_1 \times \mathbf{F}_1 + \mathbf{r}_2 \times \mathbf{F}_2 + \dots + \mathbf{r}_n \times \mathbf{F}_n
\]
где $\mathbf{r}_i$ --- радиус-вектор из точки $O$ к точке приложения $i$-й силы.

Для моментов относительно одной и той же оси вращения сложение также является векторным:
\[
\mathbf{M}_{\text{ось}} = \sum_{i=1}^n \mathbf{M}_{i,\text{ось}}
\]

\subsubsection{Момент пары сил}
\textbf{Пара сил} --- система двух равных по модулю, параллельных и противоположно направленных сил ($\mathbf{F}_1 = -\mathbf{F}_2$), не лежащих на одной прямой.

Суммарный момент пары сил \textit{не зависит} от выбора точки О и равен произведению модуля одной из сил на расстояние между линиями их действия $l$ (плечо пары):
\[
M_{пары} = F \cdot l
\]
В отличие от произвольной силы, пара сил вызывает только вращение без ускорения центра масс (так как $\sum \mathbf{F} = 0$).

\subsection{Виды равновесия}
Тело находится в равновесии, если $\sum \mathbf{F} = 0$ и $\sum \mathbf{M} = 0$. Тип равновесия определяется поведением тела при малом отклонении:

\begin{enumerate}
    \item \textbf{Устойчивое:} При отклонении возникают силы/моменты, возвращающие тело в исходное состояние. Соответствует \textbf{минимуму} потенциальной энергии ($U = \min$).
    \item \textbf{Неустойчивое:} При отклонении возникают силы, удаляющие тело от положения равновесия. Соответствует \textbf{максимуму} потенциальной энергии ($U = \max$).
    \item \textbf{Безразличное:} При отклонении тело остается в новом положении. Потенциальная энергия постоянна ($U = \text{const}$).
\end{enumerate}

\subsection{Момент импульса}

\subsubsection{Определение}

\textbf{Момент импульса} --- динамическая характеристика вращательного движения.

\begin{center}
\begin{tblr}{
    colspec = {X[c]X[l]},
    vlines,
    hlines,
    row{1} = {font=\bfseries, mode=text},
    row{odd} = {bg=gray!10},
    cells = {m},
    rowsep = 6pt
}
Формула & Название \\
$\displaystyle \mathbf{L} = [\mathbf{r}, \mathbf{p}] = m[\mathbf{r}, \mathbf{v}]$ & \textit{Момент импульса материальной точки:} векторное произведение радиус-вектора $\mathbf{r}$ на импульс $\mathbf{p}$ \\
$\displaystyle L_z = I_z \omega$ & \textit{Момент импульса твёрдого тела:} относительно неподвижной оси вращения; $I_z$ --- момент инерции, $\omega$ --- угловая скорость \\
\end{tblr}
\end{center}

\subsubsection{Основное уравнение динамики}

Рассмотрим материальную точку массой $m$, положение которой задается радиус-вектором $\mathbf{r}$ относительно начала координат $O$.

Найдем скорость изменения момента импульса:
\[
\frac{d\mathbf{L}}{dt} = \frac{d}{dt} [\mathbf{r}, m\mathbf{v}] = \left[\frac{d\mathbf{r}}{dt}, m\mathbf{v}\right] + \left[\mathbf{r}, m\frac{d\mathbf{v}}{dt}\right]
\]

В первом слагаемом $\frac{d\mathbf{r}}{dt} = \mathbf{v}$. Тогда имеем $[\mathbf{v}, m\mathbf{v}] = m[\mathbf{v}, \mathbf{v}]$. 
    Поскольку векторное произведение вектора на самого себя равно нулю, то:
    \[
    [\mathbf{v}, m\mathbf{v}] = 0
    \]
Во втором слагаемом $m\frac{d\mathbf{v}}{dt} = m\mathbf{a} = \mathbf{F}$ \textit{(по II закону Ньютона)}:
    \[
    [\mathbf{r}, \mathbf{F}] = \mathbf{M}
    \]
    где $\mathbf{M}$ --- момент силы относительно точки $O$.

Таким образом:
\[\boxed{
\frac{d\mathbf{L}}{dt} = \mathbf{M}
}\]

\subsection{Закон сохранения момента импульса}
Из основного уравнения динамики следует, что если суммарный момент внешних сил, действующих на систему, равен нулю, то момент импульса системы остается постоянным:
\[
\mathbf{L} = \text{const} \implies I_1 \omega_1 = I_2 \omega_2
\]